% GENERATED FILE. Edit canonical lessons, then run npm run generate:books.
% canonical-source-hash: fnv1a64:f0c1cd1860056c7f
% canonical-lessons: GU-C56-gaavu, GU-C56-naachvu, GU-C56-hasvu, GU-C56-saambhalvu, GU-C56-shikhvu

\chapter{Sing, Dance, Laugh, Listen, Learn}
\label{ch:gu-sing-dance-laugh-listen-learn}
\hlchaptermodality{56}

\begin{chapteropening}
\textbf{By the end of this chapter:} \emph{I can say to sing, to dance, to laugh, to listen and to learn.}

Complete the last lesson of chapter 56: I can say to sing, to dance, to laugh, to listen and to learn.
\end{chapteropening}

\section[\texorpdfstring{\gu{ગાવું} (gāvũ)}{ગાવું (gāvũ)}]{\gu{ગાવું} (gāvũ) — to sing}

\label{lesson:GU-C56-gaavu}



\begin{quote}
Before the new one: say the Gujarati for to bring, then the Gujarati for to drink.
\end{quote}

\subsection*{You'll want to know: \gu{ગાવું}}
\textbf{\gu{ગાવું}} (\emph{gāvũ}) — \textquotedblleft{}to sing\textquotedblright{}.

\begin{cousinweb}
\textbf{\gu{ગાવું}} (\emph{gāvũ}) is a tadbhava, a word inherited from Sanskrit and reshaped by sound change. It comes from \emph{gāyati}, \textquotedblleft{}sings\textquotedblright{} (CDIAL 4135).
\end{cousinweb}

\begin{grammarlens}[title={What you've built}]
One more word for sing, dance, laugh, listen, learn.
\end{grammarlens}

\subsection*{Your turn}
\begin{itemize}
\raggedright
  \item \emph{Say it:} \emph{gāvũ}
  \item \emph{Say it:} \emph{gāvũ}, once more
  \item \emph{Say it:} say \emph{gāvũ}
  \item \emph{Recall:} say the Gujarati for to bring, then the Gujarati for to drink, then say \emph{gāvũ} again
\end{itemize}

\begin{tcolorbox}[breakable,colback=teal!4,colframe=teal!35!black,title={Before you move on}]
What does \gu{ગાવું} mean? (\textquotedblleft{}To sing\textquotedblright{}.) Say it once more.
\end{tcolorbox}
\section[\texorpdfstring{\gu{નાચવું} (nācvũ)}{નાચવું (nācvũ)}]{\gu{નાચવું} (nācvũ) — to dance}

\label{lesson:GU-C56-naachvu}



\begin{quote}
Before the new one: say the Gujarati for to drink, then the Gujarati for to sing.
\end{quote}

\subsection*{You'll want to know: \gu{નાચવું}}
\textbf{\gu{નાચવું}} (\emph{nācvũ}) — \textquotedblleft{}to dance\textquotedblright{}.

\begin{cousinweb}
\textbf{\gu{નાચવું}} (\emph{nācvũ}) is a tadbhava, a word inherited from Sanskrit and reshaped by sound change. It comes from \emph{nṛtyati}, \textquotedblleft{}dances\textquotedblright{} (CDIAL 7583).
\end{cousinweb}

\begin{grammarlens}[title={What you've built}]
One more word for sing, dance, laugh, listen, learn.
\end{grammarlens}

\subsection*{Your turn}
\begin{itemize}
\raggedright
  \item \emph{Say it:} \emph{nācvũ}
  \item \emph{Say it:} \emph{nācvũ}, once more
  \item \emph{Say it:} say \emph{nācvũ}
  \item \emph{Recall:} say the Gujarati for to drink, then the Gujarati for to sing, then say \emph{nācvũ} again
\end{itemize}

\begin{tcolorbox}[breakable,colback=teal!4,colframe=teal!35!black,title={Before you move on}]
What does \gu{નાચવું} mean? (\textquotedblleft{}To dance\textquotedblright{}.) Say it once more.
\end{tcolorbox}
\section[\texorpdfstring{\gu{હસવું} (hasvũ)}{હસવું (hasvũ)}]{\gu{હસવું} (hasvũ) — to laugh}

\label{lesson:GU-C56-hasvu}



\begin{quote}
Before the new one: say the Gujarati for to sing, then the Gujarati for to dance.
\end{quote}

\subsection*{You'll want to know: \gu{હસવું}}
\textbf{\gu{હસવું}} (\emph{hasvũ}) — \textquotedblleft{}to laugh\textquotedblright{}.

\begin{cousinweb}
\textbf{\gu{હસવું}} (\emph{hasvũ}) is a tadbhava, a word inherited from Sanskrit and reshaped by sound change. It comes from \emph{hasati}, \textquotedblleft{}laughs\textquotedblright{} (CDIAL 14021).
\end{cousinweb}

\begin{grammarlens}[title={What you've built}]
One more word for sing, dance, laugh, listen, learn.
\end{grammarlens}

\subsection*{Your turn}
\begin{itemize}
\raggedright
  \item \emph{Say it:} \emph{hasvũ}
  \item \emph{Say it:} \emph{hasvũ}, once more
  \item \emph{Say it:} say \emph{hasvũ}
  \item \emph{Recall:} say the Gujarati for to sing, then the Gujarati for to dance, then say \emph{hasvũ} again
\end{itemize}

\begin{tcolorbox}[breakable,colback=teal!4,colframe=teal!35!black,title={Before you move on}]
What does \gu{હસવું} mean? (\textquotedblleft{}To laugh\textquotedblright{}.) Say it once more.
\end{tcolorbox}
\section[\texorpdfstring{\gu{સાંભળવું} (sā̃bhaḷvũ)}{સાંભળવું (sā̃bhaḷvũ)}]{\gu{સાંભળવું} (sā̃bhaḷvũ) — to listen, to hear}

\label{lesson:GU-C56-saambhalvu}



\begin{quote}
Before the new one: say the Gujarati for to dance, then the Gujarati for to laugh.
\end{quote}

\subsection*{You'll want to know: \gu{સાંભળવું}}
\textbf{\gu{સાંભળવું}} (\emph{sā̃bhaḷvũ}) — \textquotedblleft{}to listen, to hear\textquotedblright{}.

\begin{cousinweb}
\textbf{\gu{સાંભળવું}} (\emph{sāṁbhaḷvũ}) is a tadbhava, a word inherited and reshaped by sound change. It comes from \emph{*saṁbhalati}, \textquotedblleft{}attends to\textquotedblright{}, a form CDIAL reconstructs from its descendants (CDIAL 12962).
\end{cousinweb}

\begin{grammarlens}[title={What you've built}]
One more word for sing, dance, laugh, listen, learn.
\end{grammarlens}

\subsection*{Your turn}
\begin{itemize}
\raggedright
  \item \emph{Say it:} \emph{sā̃bhaḷvũ}
  \item \emph{Say it:} \emph{sā̃bhaḷvũ}, once more
  \item \emph{Say it:} say \emph{sā̃bhaḷvũ}
  \item \emph{Recall:} say the Gujarati for to dance, then the Gujarati for to laugh, then say \emph{sā̃bhaḷvũ} again
\end{itemize}

\begin{tcolorbox}[breakable,colback=teal!4,colframe=teal!35!black,title={Before you move on}]
What does \gu{સાંભળવું} mean? (\textquotedblleft{}To listen, to hear\textquotedblright{}.) Say it once more.
\end{tcolorbox}
\section[\texorpdfstring{\gu{શીખવું} (śīkhvũ)}{શીખવું (śīkhvũ)}]{\gu{શીખવું} (śīkhvũ) — to learn}

\label{lesson:GU-C56-shikhvu}



\begin{quote}
Before the new one: say the Gujarati for to laugh, then the Gujarati for to listen.
\end{quote}

\subsection*{You'll want to know: \gu{શીખવું}}
\textbf{\gu{શીખવું}} (\emph{śīkhvũ}) — \textquotedblleft{}to learn\textquotedblright{}. What you are doing now.

\begin{cousinweb}
\textbf{\gu{શીખવું}} (\emph{śīkhvũ}) is a tadbhava, a word inherited from Sanskrit and reshaped by sound change. It comes from \emph{śikṣate}, \textquotedblleft{}learns\textquotedblright{} (CDIAL 12430).
\end{cousinweb}

\begin{grammarlens}[title={What you've built}]
One more word for sing, dance, laugh, listen, learn.
\end{grammarlens}

\subsection*{Your turn}
\begin{itemize}
\raggedright
  \item \emph{Say it:} \emph{śīkhvũ}
  \item \emph{Say it:} \emph{śīkhvũ}, once more
  \item \emph{Say it:} say \emph{śīkhvũ}
  \item \emph{Recall:} say the Gujarati for to laugh, then the Gujarati for to listen, then say \emph{śīkhvũ} again
\end{itemize}

\begin{tcolorbox}[breakable,colback=teal!4,colframe=teal!35!black,title={Before you move on}]
What does \gu{શીખવું} mean? (\textquotedblleft{}To learn\textquotedblright{}.) Say it once more.
\end{tcolorbox}
